% !PRINTED 23
\weberpage{23}
\thispagestyle{empty}
\begin{center}
{\large\gesperrt{First Section.}}

\vspace{0.8em}
{\Large\gesperrt{Rational Functions.}}

\vspace{0.8em}
\rule{0.16\linewidth}{0.4pt}

\vspace{0.8em}
{\large \S.\ 1.}

\vspace{0.5em}
{\large\gesperrt{Whole Functions.}}
\end{center}

The next object of consideration is \gesperrt{whole rational}
functions, or briefly \gesperrt{whole functions of one variable}. By
these we understand expressions of the following form:
\begin{equation}
 f(x)=a_0x^n+a_1x^{n-1}+a_2x^{n-2}+\cdots+a_{n-1}x+a_n,
\tag{1}
\end{equation}
where $n$ is an integral exponent, the \gesperrt{degree} of the function $f(x)$.
The degree is a \gesperrt{natural number}. Sometimes, however, it is
also useful to speak of whole rational functions of degree 0, by
which a constant is understood. $x$ is called the \gesperrt{variable},
$a_0,a_1,a_2,\ldots,a_{n-1},a_n$ the \gesperrt{coefficients}. Both $x$
and $a_0,a_1,\ldots,a_n$ are symbols for indeterminate quantities, with
which one proceeds according to the rules of literal calculation, and
for which definite numerical values may also be substituted under
some circumstances (cf. the Introduction).

If the function $f(x)$ is written in the way shown in (1), we call it
\gesperrt{ordered according to descending powers of $x$}. The summands
may be arranged in any other order; for example, the function may also
be ordered according to ascending powers of $x$:
\[
 f(x)=a_n+a_{n-1}x+\cdots+a_1x^{n-1}+a_0x^n.
\]
By addition (subtraction) and multiplication of whole rational
functions there again arise whole rational functions. The rules of
literal calculation directly give the laws of formation of these new
functions.

In addition, the coefficient of any power $x^\nu$ in the sum is equal
to the sum of the coefficients of $x^\nu$ in the individual summands.
The degree of the function that arises is
%CONT>>
